%====================================================
%                    PRUEBAS DE CAJA BLANCA
%====================================================

\section{Pruebas de Caja Blanca}

%====================================================
%               CU-01  ---> IniciarSesion
%====================================================

\subsection{Caso de prueba: CU01}

\begin{table}[H]
\centering
\label{tab:cu01_white_test}
\renewcommand{\arraystretch}{1.3}
\begin{tabular}{|c|p{5cm}|p{6cm}|p{2.5cm}|}
\hline
\rowcolor[RGB]{217,217,217}
\textbf{Caso} & \textbf{Descripción} & \textbf{Resultado esperado} & \textbf{Resultado obtenido} \\
\hline

1 &
El usuario intenta iniciar sesión con credenciales válidas. &
Se autentica correctamente y se generan tokens JWT. &
 \\
\hline

2 &
El usuario intenta iniciar sesión con contraseña incorrecta. &
Se rechaza la autenticación y se retorna error 401. &
 \\
\hline

3 &
El usuario intenta iniciar sesión con email no registrado. &
Se retorna error 404 indicando usuario no encontrado. &
 \\
\hline

\end{tabular}
\caption{Caso de prueba: CU01.}
\end{table}

%====================================================
%               CU-02  ---> Registro
%====================================================

\subsection{Caso de prueba: CU02}

\begin{table}[H]
\centering
\label{tab:cu02_white_test}
\renewcommand{\arraystretch}{1.3}
\begin{tabular}{|c|p{5cm}|p{6cm}|p{2.5cm}|}
\hline
\rowcolor[RGB]{217,217,217}
\textbf{Caso} & \textbf{Descripción} & \textbf{Resultado esperado} & \textbf{Resultado obtenido} \\
\hline

1 &
El usuario intenta registrarse con datos válidos. &
Se crea la cuenta en la BD y se retorna el nuevo usuario. &
 \\
\hline

2 &
El usuario intenta registrarse con email duplicado. &
Se rechaza la operación y se retorna error 409. &
 \\
\hline

3 &
El usuario intenta registrarse sin completar campos requeridos. &
Se validan los datos y se retorna error 400. &
 \\
\hline

\end{tabular}
\caption{Caso de prueba: CU02.}
\end{table}

%====================================================
%               CU-03  ---> CerrarSesion
%====================================================

\subsection{Caso de prueba: CU03}

\begin{table}[H]
\centering
\label{tab:cu03_white_test}
\renewcommand{\arraystretch}{1.3}
\begin{tabular}{|c|p{5cm}|p{6cm}|p{2.5cm}|}
\hline
\rowcolor[RGB]{217,217,217}
\textbf{Caso} & \textbf{Descripción} & \textbf{Resultado esperado} & \textbf{Resultado obtenido} \\
\hline

1 &
El usuario cierra sesión con token válido. &
Se invalida el token y se retorna confirmación. &
 \\
\hline

2 &
El usuario intenta cerrar sesión sin token. &
Se retorna error 401 no autorizado. &
 \\
\hline

\end{tabular}
\caption{Caso de prueba: CU03.}
\end{table}

%====================================================
%               CU-04  ---> VerPerfil
%====================================================

\subsection{Caso de prueba: CU04}

\begin{table}[H]
\centering
\label{tab:cu04_white_test}
\renewcommand{\arraystretch}{1.3}
\begin{tabular}{|c|p{5cm}|p{6cm}|p{2.5cm}|}
\hline
\rowcolor[RGB]{217,217,217}
\textbf{Caso} & \textbf{Descripción} & \textbf{Resultado esperado} & \textbf{Resultado obtenido} \\
\hline

1 &
El usuario accede a su perfil con token válido. &
Se retornan los datos del perfil del usuario. &
 \\
\hline

2 &
El usuario intenta acceder al perfil sin autenticación. &
Se retorna error 401. &
 \\
\hline

\end{tabular}
\caption{Caso de prueba: CU04.}
\end{table}

%====================================================
%               CU-05  ---> ModificarPerfil
%====================================================

\subsection{Caso de prueba: CU05}

\begin{table}[H]
\centering
\label{tab:cu05_white_test}
\renewcommand{\arraystretch}{1.3}
\begin{tabular}{|c|p{5cm}|p{6cm}|p{2.5cm}|}
\hline
\rowcolor[RGB]{217,217,217}
\textbf{Caso} & \textbf{Descripción} & \textbf{Resultado esperado} & \textbf{Resultado obtenido} \\
\hline

1 &
El usuario modifica su perfil con datos válidos. &
Los cambios se guardan en la BD y se retorna el perfil actualizado. &
 \\
\hline

2 &
El usuario intenta usar un email ya registrado. &
Se rechaza la operación y se retorna error 409. &
 \\
\hline

3 &
El usuario intenta modificar el perfil sin autenticación. &
Se retorna error 401. &
 \\
\hline

\end{tabular}
\caption{Caso de prueba: CU05.}
\end{table}

%====================================================
%               CU-06  ---> EliminarCuenta
%====================================================

\subsection{Caso de prueba: CU06}

\begin{table}[H]
\centering
\label{tab:cu06_white_test}
\renewcommand{\arraystretch}{1.3}
\begin{tabular}{|c|p{5cm}|p{6cm}|p{2.5cm}|}
\hline
\rowcolor[RGB]{217,217,217}
\textbf{Caso} & \textbf{Descripción} & \textbf{Resultado esperado} & \textbf{Resultado obtenido} \\
\hline

1 &
El usuario elimina su cuenta con contraseña correcta. &
Se elimina el usuario y todos sus datos asociados. &
 \\
\hline

2 &
El usuario intenta eliminar la cuenta con contraseña incorrecta. &
Se retorna error 401. &
 \\
\hline

3 &
El usuario intenta eliminar sin autenticación. &
Se retorna error 401. &
 \\
\hline

\end{tabular}
\caption{Caso de prueba: CU06.}
\end{table}

%====================================================
%               CU-07  ---> RealizarReserva
%====================================================

\subsection{Caso de prueba: CU07}

\begin{table}[H]
\centering
\label{tab:cu07_white_test}
\renewcommand{\arraystretch}{1.3}
\begin{tabular}{|c|p{5cm}|p{6cm}|p{2.5cm}|}
\hline
\rowcolor[RGB]{217,217,217}
\textbf{Caso} & \textbf{Descripción} & \textbf{Resultado esperado} & \textbf{Resultado obtenido} \\
\hline

1 &
El usuario realiza una reserva con datos válidos. &
Se crea la reserva en la BD y se envía notificación. &
 \\
\hline

2 &
El usuario intenta reservar una hora no disponible. &
Se retorna error 409 indicando conflicto de horario. &
 \\
\hline

3 &
El usuario intenta hacer reserva sin autenticación. &
Se retorna error 401. &
 \\
\hline

\end{tabular}
\caption{Caso de prueba: CU07.}
\end{table}

%====================================================
%               CU-08  ---> VerReservas
%====================================================

\subsection{Caso de prueba: CU08}

\begin{table}[H]
\centering
\label{tab:cu08_white_test}
\renewcommand{\arraystretch}{1.3}
\begin{tabular}{|c|p{5cm}|p{6cm}|p{2.5cm}|}
\hline
\rowcolor[RGB]{217,217,217}
\textbf{Caso} & \textbf{Descripción} & \textbf{Resultado esperado} & \textbf{Resultado obtenido} \\
\hline

1 &
El usuario consulta sus reservas con filtros válidos. &
Se retorna lista de reservas del usuario. &
 \\
\hline

2 &
El usuario intenta ver reservas sin autenticación. &
Se retorna error 401. &
 \\
\hline

\end{tabular}
\caption{Caso de prueba: CU08.}
\end{table}

%====================================================
%               CU-09  ---> CancelarReserva
%====================================================

\subsection{Caso de prueba: CU09}

\begin{table}[H]
\centering
\label{tab:cu09_white_test}
\renewcommand{\arraystretch}{1.3}
\begin{tabular}{|c|p{5cm}|p{6cm}|p{2.5cm}|}
\hline
\rowcolor[RGB]{217,217,217}
\textbf{Caso} & \textbf{Descripción} & \textbf{Resultado esperado} & \textbf{Resultado obtenido} \\
\hline

1 &
El usuario cancela una reserva válida. &
Se actualiza el estado y se envía notificación. &
 \\
\hline

2 &
El usuario intenta cancelar una reserva inexistente. &
Se retorna error 404. &
 \\
\hline

3 &
El usuario intenta cancelar sin autenticación. &
Se retorna error 401. &
 \\
\hline

\end{tabular}
\caption{Caso de prueba: CU09.}
\end{table}

%====================================================
%               CU-10  ---> ExportarEventos
%====================================================

\subsection{Caso de prueba: CU10}

\begin{table}[H]
\centering
\label{tab:cu10_white_test}
\renewcommand{\arraystretch}{1.3}
\begin{tabular}{|c|p{5cm}|p{6cm}|p{2.5cm}|}
\hline
\rowcolor[RGB]{217,217,217}
\textbf{Caso} & \textbf{Descripción} & \textbf{Resultado esperado} & \textbf{Resultado obtenido} \\
\hline

1 &
El usuario exporta sus reservas a Google Calendar. &
Los eventos se crean en Google Calendar del usuario. &
 \\
\hline

2 &
El usuario intenta exportar sin autenticación OAuth. &
Se solicita autenticación. &
 \\
\hline

\end{tabular}
\caption{Caso de prueba: CU10.}
\end{table}

%====================================================
%               CU-11  ---> VerNotificaciones
%====================================================

\subsection{Caso de prueba: CU11}

\begin{table}[H]
\centering
\label{tab:cu11_white_test}
\renewcommand{\arraystretch}{1.3}
\begin{tabular}{|c|p{5cm}|p{6cm}|p{2.5cm}|}
\hline
\rowcolor[RGB]{217,217,217}
\textbf{Caso} & \textbf{Descripción} & \textbf{Resultado esperado} & \textbf{Resultado obtenido} \\
\hline

1 &
El usuario consulta su historial de notificaciones. &
Se retorna lista de notificaciones. &
 \\
\hline

2 &
El usuario intenta ver notificaciones sin autenticación. &
Se retorna error 401. &
 \\
\hline

\end{tabular}
\caption{Caso de prueba: CU11.}
\end{table}

%====================================================
%               CU-12  ---> LeerNotificacion
%====================================================

\subsection{Caso de prueba: CU12}

\begin{table}[H]
\centering
\label{tab:cu12_white_test}
\renewcommand{\arraystretch}{1.3}
\begin{tabular}{|c|p{5cm}|p{6cm}|p{2.5cm}|}
\hline
\rowcolor[RGB]{217,217,217}
\textbf{Caso} & \textbf{Descripción} & \textbf{Resultado esperado} & \textbf{Resultado obtenido} \\
\hline

1 &
El usuario marca una notificación como leída. &
Se actualiza el estado en la BD. &
 \\
\hline

2 &
El usuario intenta marcar notificación sin autenticación. &
Se retorna error 401. &
 \\
\hline

\end{tabular}
\caption{Caso de prueba: CU12.}
\end{table}

%====================================================
%               CU-13  ---> RecalcularNegocios
%====================================================

\subsection{Caso de prueba: CU13}

\begin{table}[H]
\centering
\label{tab:cu13_white_test}
\renewcommand{\arraystretch}{1.3}
\begin{tabular}{|c|p{5cm}|p{6cm}|p{2.5cm}|}
\hline
\rowcolor[RGB]{217,217,217}
\textbf{Caso} & \textbf{Descripción} & \textbf{Resultado esperado} & \textbf{Resultado obtenido} \\
\hline

1 &
Se recalculan los negocios después de cambios. &
Los datos se actualizan correctamente en la BD. &
 \\
\hline

\end{tabular}
\caption{Caso de prueba: CU13.}
\end{table}

%====================================================
%               CU-14  ---> AplicarFiltros
%====================================================

\subsection{Caso de prueba: CU14}

\begin{table}[H]
\centering
\label{tab:cu14_white_test}
\renewcommand{\arraystretch}{1.3}
\begin{tabular}{|c|p{5cm}|p{6cm}|p{2.5cm}|}
\hline
\rowcolor[RGB]{217,217,217}
\textbf{Caso} & \textbf{Descripción} & \textbf{Resultado esperado} & \textbf{Resultado obtenido} \\
\hline

1 &
El usuario aplica filtros a la búsqueda de negocios. &
Se retornan negocios que coinciden con los criterios. &
 \\
\hline

2 &
El usuario aplica filtros sin criterios válidos. &
Se retorna error 400. &
 \\
\hline

\end{tabular}
\caption{Caso de prueba: CU14.}
\end{table}

%====================================================
%               CU-15  ---> BuscarNegocio
%====================================================

\subsection{Caso de prueba: CU15}

\begin{table}[H]
\centering
\label{tab:cu15_white_test}
\renewcommand{\arraystretch}{1.3}
\begin{tabular}{|c|p{5cm}|p{6cm}|p{2.5cm}|}
\hline
\rowcolor[RGB]{217,217,217}
\textbf{Caso} & \textbf{Descripción} & \textbf{Resultado esperado} & \textbf{Resultado obtenido} \\
\hline

1 &
El usuario busca negocios por nombre o categoría. &
Se retornan negocios que coinciden. &
 \\
\hline

2 &
El usuario realiza búsqueda con término vacío. &
Se retornan todos los negocios o error 400. &
 \\
\hline

\end{tabular}
\caption{Caso de prueba: CU15.}
\end{table}

%====================================================
%               CU-16  ---> OrientarMapa
%====================================================

\subsection{Caso de prueba: CU16}

\begin{table}[H]
\centering
\label{tab:cu16_white_test}
\renewcommand{\arraystretch}{1.3}
\begin{tabular}{|c|p{5cm}|p{6cm}|p{2.5cm}|}
\hline
\rowcolor[RGB]{217,217,217}
\textbf{Caso} & \textbf{Descripción} & \textbf{Resultado esperado} & \textbf{Resultado obtenido} \\
\hline

1 &
El usuario obtiene su ubicación actual en el mapa. &
Se muestra el mapa centrado en su ubicación. &
 \\
\hline

2 &
El usuario rechaza permisos de localización. &
Se muestra error o ubicación por defecto. &
 \\
\hline

\end{tabular}
\caption{Caso de prueba: CU16.}
\end{table}

%====================================================
%               CU-17  ---> CalcularRuta
%====================================================

\subsection{Caso de prueba: CU17}

\begin{table}[H]
\centering
\label{tab:cu17_white_test}
\renewcommand{\arraystretch}{1.3}
\begin{tabular}{|c|p{5cm}|p{6cm}|p{2.5cm}|}
\hline
\rowcolor[RGB]{217,217,217}
\textbf{Caso} & \textbf{Descripción} & \textbf{Resultado esperado} & \textbf{Resultado obtenido} \\
\hline

1 &
El usuario calcula ruta a un negocio. &
Se muestra la ruta más óptima en el mapa. &
 \\
\hline

2 &
No hay ruta disponible hacia el destino. &
Se retorna error indicando problema de ruta. &
 \\
\hline

\end{tabular}
\caption{Caso de prueba: CU17.}
\end{table}

%====================================================
%               CU-18  ---> VerDetallesNegocio
%====================================================

\subsection{Caso de prueba: CU18}

\begin{table}[H]
\centering
\label{tab:cu18_white_test}
\renewcommand{\arraystretch}{1.3}
\begin{tabular}{|c|p{5cm}|p{6cm}|p{2.5cm}|}
\hline
\rowcolor[RGB]{217,217,217}
\textbf{Caso} & \textbf{Descripción} & \textbf{Resultado esperado} & \textbf{Resultado obtenido} \\
\hline

1 &
El usuario visualiza detalles de un negocio. &
Se retornan todos los datos del negocio. &
 \\
\hline

2 &
El usuario intenta ver detalles de negocio inexistente. &
Se retorna error 404. &
 \\
\hline

\end{tabular}
\caption{Caso de prueba: CU18.}
\end{table}

%====================================================
%               CU-19  ---> VincularNegocio
%====================================================

\subsection{Caso de prueba: CU19}

\begin{table}[H]
\centering
\label{tab:cu19_white_test}
\renewcommand{\arraystretch}{1.3}
\begin{tabular}{|c|p{5cm}|p{6cm}|p{2.5cm}|}
\hline
\rowcolor[RGB]{217,217,217}
\textbf{Caso} & \textbf{Descripción} & \textbf{Resultado esperado} & \textbf{Resultado obtenido} \\
\hline

1 &
El propietario vincula su negocio con Google Places. &
Se asocian los datos correctamente. &
 \\
\hline

2 &
Se intenta vincular un negocio ya vinculado. &
Se retorna error 409. &
 \\
\hline

3 &
Se intenta vincular sin autenticación. &
Se retorna error 401. &
 \\
\hline

\end{tabular}
\caption{Caso de prueba: CU19.}
\end{table}

%====================================================
%               CU-20  ---> EditarNegocio
%====================================================

\subsection{Caso de prueba: CU20}

\begin{table}[H]
\centering
\label{tab:cu20_white_test}
\renewcommand{\arraystretch}{1.3}
\begin{tabular}{|c|p{5cm}|p{6cm}|p{2.5cm}|}
\hline
\rowcolor[RGB]{217,217,217}
\textbf{Caso} & \textbf{Descripción} & \textbf{Resultado esperado} & \textbf{Resultado obtenido} \\
\hline

1 &
El propietario edita datos de su negocio. &
Los cambios se guardan en la BD. &
 \\
\hline

2 &
El propietario intenta editar datos inválidos. &
Se retorna error 400. &
 \\
\hline

3 &
Se intenta editar sin ser el propietario. &
Se retorna error 403 prohibido. &
 \\
\hline

\end{tabular}
\caption{Caso de prueba: CU20.}
\end{table}

%====================================================
%               CU-21  ---> VerEstadisticas
%====================================================

\subsection{Caso de prueba: CU21}

\begin{table}[H]
\centering
\label{tab:cu21_white_test}
\renewcommand{\arraystretch}{1.3}
\begin{tabular}{|c|p{5cm}|p{6cm}|p{2.5cm}|}
\hline
\rowcolor[RGB]{217,217,217}
\textbf{Caso} & \textbf{Descripción} & \textbf{Resultado esperado} & \textbf{Resultado obtenido} \\
\hline

1 &
El propietario visualiza estadísticas de su negocio. &
Se retornan datos de reservas, ingresos, etc. &
 \\
\hline

2 &
Se intenta consultar estadísticas sin autenticación. &
Se retorna error 401. &
 \\
\hline

\end{tabular}
\caption{Caso de prueba: CU21.}
\end{table}

%====================================================
%               CU-22  ---> MarcarVacaciones
%====================================================

En la Tabla \ref{tab:cu22_white_test} se muestran las pruebas realizadas sobre el Caso de Uso: MarcarVacaciones.

\subsection{Caso de prueba: CU22}

\begin{table}[H]
\centering
\label{tab:cu22_white_test}
\renewcommand{\arraystretch}{1.3}
\begin{tabular}{|c|p{5cm}|p{6cm}|p{2.5cm}|}
\hline
\rowcolor[RGB]{217,217,217}
\textbf{Caso} & \textbf{Descripción} & \textbf{Resultado esperado} & \textbf{Resultado obtenido} \\
\hline

1 &
El propietario marca un rango de fechas como vacaciones. &
Se registran las fechas como períodos indisponibles en la BD. &
 \\
\hline

2 &
El propietario intenta marcar fechas pasadas como vacaciones. &
Se rechaza la operación y se retorna error de fecha inválida. &
 \\
\hline

3 &
El propietario marca vacaciones sin tener un negocio vinculado. &
Se retorna error 400 indicando que debe tener un negocio activo. &
 \\
\hline

\end{tabular}
\caption{Caso de prueba: CU22.}
\label{tab:cu22_white_test}
\end{table}

%====================================================
%               CU-23  ---> DesvincularNegocio
%====================================================

En la Tabla \ref{tab:cu23_white_test} se muestran las pruebas realizadas sobre el Caso de Uso: DesvincularNegocio.

\subsection{Caso de prueba: CU23}

\begin{table}[H]
\centering
\label{tab:cu23_white_test}
\renewcommand{\arraystretch}{1.3}
\begin{tabular}{|c|p{5cm}|p{6cm}|p{2.5cm}|}
\hline
\rowcolor[RGB]{217,217,217}
\textbf{Caso} & \textbf{Descripción} & \textbf{Resultado esperado} & \textbf{Resultado obtenido} \\
\hline

1 &
El propietario desvincula un negocio de su cuenta. &
Se elimina la asociación del negocio a la cuenta. &
 \\
\hline

2 &
El propietario intenta desvincularse de un negocio no existente. &
Se retorna error 404 indicando que el negocio no existe. &
 \\
\hline

3 &
Se intenta desvincularse sin tener permisos de propietario. &
Se retorna error 403 (Forbidden). &
 \\
\hline

\end{tabular}
\caption{Caso de prueba: CU23.}
\label{tab:cu23_white_test}
\end{table}

%====================================================
%               CU-24  ---> RealizarCitas
%====================================================

En la Tabla \ref{tab:cu24_white_test} se muestran las pruebas realizadas sobre el Caso de Uso: RealizarCitas.

\subsection{Caso de prueba: CU24}

\begin{table}[H]
\centering
\label{tab:cu24_white_test}
\renewcommand{\arraystretch}{1.3}
\begin{tabular}{|c|p{5cm}|p{6cm}|p{2.5cm}|}
\hline
\rowcolor[RGB]{217,217,217}
\textbf{Caso} & \textbf{Descripción} & \textbf{Resultado esperado} & \textbf{Resultado obtenido} \\
\hline

1 &
El propietario agenda una cita manualmente con datos válidos. &
Se registra la cita en la BD y se retorna confirmación. &
 \\
\hline

2 &
El propietario intenta agendar una cita en horario no disponible. &
Se rechaza la operación y retorna error de conflicto. &
 \\
\hline

3 &
El propietario agenda una cita sin tener un negocio activo. &
Se retorna error 400 indicando que debe seleccionar un negocio. &
 \\
\hline

\end{tabular}
\caption{Caso de prueba: CU24.}
\label{tab:cu24_white_test}
\end{table}

%====================================================
%               CU-25  ---> VerCitas
%====================================================

En la Tabla \ref{tab:cu25_white_test} se muestran las pruebas realizadas sobre el Caso de Uso: VerCitas.

\subsection{Caso de prueba: CU25}

\begin{table}[H]
\centering
\label{tab:cu25_white_test}
\renewcommand{\arraystretch}{1.3}
\begin{tabular}{|c|p{5cm}|p{6cm}|p{2.5cm}|}
\hline
\rowcolor[RGB]{217,217,217}
\textbf{Caso} & \textbf{Descripción} & \textbf{Resultado esperado} & \textbf{Resultado obtenido} \\
\hline

1 &
El propietario consulta sus citas programadas. &
Se retorna lista de todas las citas del negocio. &
 \\
\hline

2 &
El propietario consulta citas sin tener ninguna agendada. &
Se retorna lista vacía. &
 \\
\hline

3 &
Se intenta consultar citas sin estar autenticado. &
Se retorna error 401 (Unauthorized). &
 \\
\hline

\end{tabular}
\caption{Caso de prueba: CU25.}
\label{tab:cu25_white_test}
\end{table}

%====================================================
%               CU-26  ---> CancelarCita
%====================================================

En la Tabla \ref{tab:cu26_white_test} se muestran las pruebas realizadas sobre el Caso de Uso: CancelarCita.

\subsection{Caso de prueba: CU26}

\begin{table}[H]
\centering
\label{tab:cu26_white_test}
\renewcommand{\arraystretch}{1.3}
\begin{tabular}{|c|p{5cm}|p{6cm}|p{2.5cm}|}
\hline
\rowcolor[RGB]{217,217,217}
\textbf{Caso} & \textbf{Descripción} & \textbf{Resultado esperado} & \textbf{Resultado obtenido} \\
\hline

1 &
El propietario cancela una cita activa. &
Se actualiza el estado de la cita a cancelada. &
 \\
\hline

2 &
El propietario intenta cancelar una cita ya cancelada. &
Se retorna error indicando que la cita ya está cancelada. &
 \\
\hline

3 &
El propietario intenta cancelar una cita inexistente. &
Se retorna error 404 (Not Found). &
 \\
\hline

\end{tabular}
\caption{Caso de prueba: CU26.}
\label{tab:cu26_white_test}
\end{table}

%====================================================
%               CU-27  ---> MarcarFavorito
%====================================================

En la Tabla \ref{tab:cu27_white_test} se muestran las pruebas realizadas sobre el Caso de Uso: MarcarFavorito.

\subsection{Caso de prueba: CU27}

\begin{table}[H]
\centering
\label{tab:cu27_white_test}
\renewcommand{\arraystretch}{1.3}
\begin{tabular}{|c|p{5cm}|p{6cm}|p{2.5cm}|}
\hline
\rowcolor[RGB]{217,217,217}
\textbf{Caso} & \textbf{Descripción} & \textbf{Resultado esperado} & \textbf{Resultado obtenido} \\
\hline

1 &
El cliente marca un negocio válido como favorito. &
Se agrega el negocio a la lista de favoritos del cliente. &
 \\
\hline

2 &
El cliente intenta marcar un negocio ya en favoritos. &
Se retorna error indicando que ya está en favoritos. &
 \\
\hline

3 &
El cliente marca un negocio no existente como favorito. &
Se retorna error 404 (Not Found). &
 \\
\hline

\end{tabular}
\caption{Caso de prueba: CU27.}
\label{tab:cu27_white_test}
\end{table}

%====================================================
%               CU-28  ---> DesmarcarFavorito
%====================================================

En la Tabla \ref{tab:cu28_white_test} se muestran las pruebas realizadas sobre el Caso de Uso: DesmarcarFavorito.

\subsection{Caso de prueba: CU28}

\begin{table}[H]
\centering
\label{tab:cu28_white_test}
\renewcommand{\arraystretch}{1.3}
\begin{tabular}{|c|p{5cm}|p{6cm}|p{2.5cm}|}
\hline
\rowcolor[RGB]{217,217,217}
\textbf{Caso} & \textbf{Descripción} & \textbf{Resultado esperado} & \textbf{Resultado obtenido} \\
\hline

1 &
El cliente desmar
ca un negocio de sus favoritos. &
Se elimina el negocio de la lista de favoritos. &
 \\
\hline

2 &
El cliente intenta desmarcar un negocio no en favoritos. &
Se retorna error indicando que no está en favoritos. &
 \\
\hline

3 &
El cliente intenta desmarcar un negocio inexistente. &
Se retorna error 404 (Not Found). &
 \\
\hline

\end{tabular}
\caption{Caso de prueba: CU28.}
\label{tab:cu28_white_test}
\end{table}
